\section{Formal Oracle Computation}

\subsection{Oracle Program Syntax}

The file \texttt{OracleComputability/OracleCode.lean} defines the program language:

\begin{lstlisting}
inductive OracleCode : Type where
  | zero
  | succ
  | left
  | right
  | query
  | pair (f g : OracleCode)
  | comp (f g : OracleCode)
  | prec (f g : OracleCode)
  | rfind (f : OracleCode)
\end{lstlisting}

The first four constructors are the usual initial functions on natural
numbers.  \texttt{pair}, \texttt{comp}, \texttt{prec}, and \texttt{rfind}
give pairing, composition, primitive recursion, and unbounded search, in
the same unary-with-pairing style used by Mathlib's
\texttt{Nat.Partrec.Code}.  The additional constructor \texttt{query}
asks the oracle whether the input number is in the oracle set, returning
\(1\) for true and \(0\) for false.

\subsection{Partial Finite Oracles}

The evaluator uses partial oracles:

\begin{lstlisting}
abbrev PartOracle : Type := Nat -> Option Bool

def PartOracle.ofFun (X : Nat -> Bool) : PartOracle :=
  fun n => some (X n)

def PartOracle.ofSnapshot (sigma : List Bool) : PartOracle :=
  fun n => sigma[n]?
\end{lstlisting}

A total characteristic function is embedded by \texttt{ofFun}; a finite
snapshot answers positions inside the list and returns \texttt{none}
outside it.  This makes insufficient oracle information a semantic case,
not an accidental timeout.

\subsection{Step-Indexed Evaluation and Results}

Successful runs return both output and use:

\begin{lstlisting}
structure HaltData where
  output : Nat
  use : Nat

inductive RunResult where
  | halt (d : HaltData)
  | stuck
  | timeout
\end{lstlisting}

The evaluator

\begin{lstlisting}
def run : Nat -> PartOracle -> OracleCode -> Nat -> RunResult
\end{lstlisting}

is total in Lean because it is indexed by fuel.  If the fuel runs out it
returns \texttt{timeout}.  If a query asks for a position where a finite
snapshot has no answer, it returns \texttt{stuck}.  If it halts, it returns
\texttt{halt \(\langle output,use\rangle\)}.  The result type distinguishes
three facts that the priority proof uses differently: convergence,
insufficient finite oracle information, and insufficient fuel.

Use is propagated through sequencing by maximum.  For a query at input
\(x\), the recorded use is \(x+1\).  For query-free initial functions the
use is \(0\).  The file \texttt{FiniteEval.lean} provides named step
lemmas such as \texttt{run_query_some}, \texttt{run_comp_step},
\texttt{run_prec_succ_step}, and \texttt{run_rfind_step}; later proofs
rewrite through these equations rather than unfolding the evaluator
directly.

\subsection{The Use Theorem}

The main preservation theorem is \texttt{run_halt_mono}:

\begin{lstlisting}
theorem run_halt_mono
    (hk : k <= k')
    (hO :
      forall i, i < d.use ->
      forall b, O i = some b -> O' i = some b)
    (h : run k O c x = .halt d) :
    run k' O' c x = .halt d
\end{lstlisting}

The theorem says that a halting bounded run survives an increase of fuel
and replacement of the oracle by one that preserves all answers actually
used by the original computation.  Notice that \(O'\) need not extend \(O\)
everywhere.  It only has to preserve answered positions below the recorded
use; positions above the use and unanswered positions are irrelevant.

The proof follows the recursion defining \texttt{run}.  In composition,
primitive recursion, and minimization, the proof decomposes a halting
\texttt{bind}, applies the induction hypothesis to subruns, and rebuilds
the result while preserving the maximum-use bookkeeping.  This is the
formal version of the classical use principle: a computation already seen
at a finite stage is preserved if future enumeration avoids its use.

\subsection{Infinite Semantics}

The file \texttt{OracleComputability/InfiniteEval.lean} defines:

\begin{lstlisting}
def Computes (c : OracleCode) (X : Nat -> Bool)
    (x y : Nat) : Prop :=
  exists k d,
    run k (PartOracle.ofFun X) c x = .halt d /\
    d.output = y
\end{lstlisting}

Thus infinite computation is recovered from existence of a finite
successful run against the total oracle.  The theorem
\texttt{Computes.unique} gives determinism of outputs.  The theorem
\texttt{computes_iff_initialSegment} connects total-oracle computation to
finite snapshots: a total computation exists iff some finite initial
segment already suffices.  This bridge is essential in the never-acted case
of the requirement proof, where a hypothetical total computation must
eventually become visible at a finite stage.

\subsection{Program Numbering}

\texttt{OracleComputability/Numbering.lean} defines an explicit coding
\texttt{encodeCode : OracleCode -> Nat} and decoding
\texttt{ofNatCode : Nat -> OracleCode}.  The two round-trip theorems
\texttt{ofNatCode_encodeCode} and \texttt{encodeCode_ofNatCode} prove that
the numbering is a bijection.  Requirements are indexed by natural numbers,
and the final diagonalization converts an arbitrary candidate program
\(c\) to its index \texttt{OracleCode.encodeCode c}; the round-trip theorem
then shows that the corresponding requirement covers exactly \(c\).

